\chapter{SVM CLASSIFIER EVALUATION}

\section{Experimental Setup}

We evaluated the Support Vector Machine (SVM) classifier across both gene-expression microarray datasets described in Chapter~3: GSE19804 (balanced lung cancer, 120 samples) and GSE42568 (imbalanced breast cancer, 121 samples). Unless otherwise stated, the SVM used a linear kernel with class weighting set to \texttt{"balanced"} and a fixed regularization strength $C = 1.0$; this is the configuration used to produce the headline Path A/Path B comparisons in this chapter (Tables~\ref{tab:gse19804_svm_results} and~\ref{tab:gse42568_svm_results}). Section~\ref{sec:imbalance_extensions} separately reports a set of tuned variants in which $C$ and the class-weight setting are instead selected by inner cross-validation; these tuned variants are evaluated only on the Welch $t$-test feature-selection path and are reported alongside, not in place of, the fixed-$C=1.0$ baseline.

The experimental framework comprised two distinct tracks:
\begin{itemize}
    \item \textbf{Path A (Baseline):} Linear SVM trained on the full high-dimensional feature space ($p = 54,675$) without feature selection.
    \item \textbf{Path B (Optimised):} Linear SVM trained on the probes retained by the selected feature-selection rule, with feature selection performed strictly \textit{inside} each training fold using \texttt{sklearn.pipeline.Pipeline}.
\end{itemize}

By placing the \texttt{FeatureSelector} inside the cross-validation loop, feature ranking and subset selection were computed exclusively on each training fold ($80\%$ of data), which is a precondition for avoiding data leakage to the evaluation test fold ($20\%$ of data). As discussed in Chapter~3, this precondition applies to every data-dependent step in the pipeline: scaling, feature selection, SMOTE, threshold tuning, and hyperparameter search.

\subsection{Evaluation Metrics}

Each pipeline is scored on five folds and reported as mean $\pm$ standard deviation. Because the two datasets differ substantially in class balance, no single metric is sufficient on its own, so the following complementary measures are reported together:

\begin{itemize}
    \item \textbf{Accuracy} : the overall proportion of correctly classified samples. On imbalanced data this metric can be misleadingly high if the classifier simply predicts the majority class.
    \item \textbf{Precision and Recall (Sensitivity)} : precision is the proportion of predicted-positive samples that are true positives; recall (equivalently, sensitivity) is the proportion of true positives that are correctly identified.
    \item \textbf{Specificity} : the proportion of true negatives correctly identified. This is the metric most sensitive to majority-class collapse, since a classifier that always predicts the positive class scores $0\%$ specificity regardless of its accuracy.
    \item \textbf{F1-Score} : the harmonic mean of precision and recall, useful when the cost of false positives and false negatives is similar.
    \item \textbf{Matthews Correlation Coefficient (MCC)} : a single balanced measure computed from all four confusion-matrix cells, ranging from $-1$ (total disagreement) through $0$ (no linear association between predicted and true labels) to $+1$ (perfect prediction). An MCC of exactly zero can arise from random guessing, but in this study it typically arises from the classifier constantly predicting the majority class, which is confirmed directly from the confusion matrix wherever it occurs. MCC is regarded as the most reliable single indicator of classifier quality under class imbalance and is treated as the primary decision metric in this chapter whenever it conflicts with accuracy.
    \item \textbf{ROC-AUC} : the area under the receiver operating characteristic curve, summarising the trade-off between sensitivity and the false-positive rate independently of the classification threshold.
\end{itemize}

\section{Imbalanced-Data Extensions and Model Tuning}
\label{sec:imbalance_extensions}

\subsection{Scope and Methodology}

The GSE42568 analysis was extended to test whether the poor minority-class performance was driven by classifier design or by class imbalance itself, though this experiment does not by itself isolate which factor is responsible (see Section~\ref{sec:gse42568_discussion}). These extensions were run for all seven Path B feature-selection methods: for each method, the following variants were evaluated, all using the same fold-local feature-selection pipeline as Table~\ref{tab:gse42568_svm_results}:

\begin{itemize}
    \item Baseline: default decision threshold (0.0), $C = 1.0$
    \item SMOTE: fold-wise resampling applied to the training fold \emph{before} feature selection
    \item Threshold tuning: Youden's $J$ statistic applied to out-of-fold decision scores (Chapter~3)
    \item $C$ tuning: inner-CV grid search for SVM regularization ($C \in \{0.01, 0.1, 1, 10, 100\}$), together with a comparison of \texttt{None} versus \texttt{"balanced"} class weighting
    \item Logistic regression, Random Forest, and XGBoost: alternative classifiers using the same fold-selected probes
\end{itemize}

\textbf{A methodological caveat on the SMOTE variant.} Because SMOTE is applied \emph{before} the feature-selection step in this pipeline, the $t$-test/ANOVA/BH-FDR filters see a synthetically class-balanced training fold rather than the true, imbalanced one. This inflates the number of probes that reach nominal significance: SMOTE retained a mean of $20{,}487$ probes under Welch selection and $13{,}259$ under BH-FDR, versus $6{,}486$ and $1.6$ respectively without SMOTE (Table~\ref{tab:gse42568_smote_inflation}). This is a known interaction in the SMOTE-plus-feature-selection literature --- synthetic samples are highly correlated with their real neighbours, which understates within-group variance and inflates apparent significance --- and is reported here as a methodological caveat rather than a genuine gain in discriminative signal, since the SMOTE variant's MCC ($-0.105$ for all three filter methods, identical because the resulting downstream classifier collapses to the majority class in exactly the same way regardless of which of the three near-identical inflated feature sets it receives) is in fact \emph{worse} than the un-resampled baseline for all three filters.

\begin{table}[H]
    \centering
    \caption{Effect of pre-selection SMOTE on the retained probe count and MCC, filter methods on GSE42568 (5-fold CV mean).}
    \label{tab:gse42568_smote_inflation}

    \small
    \resizebox{\textwidth}{!}{%
        \begin{tabular}{lcccc}
\toprule
\textbf{Method} & \textbf{Probes (no SMOTE)} & \textbf{Probes (with SMOTE)} & \textbf{MCC (no SMOTE)} & \textbf{MCC (with SMOTE)} \\
\midrule
Welch $t$-test & 6,498.0 & 20,331.0 & $-0.0315$ & $-0.1049$ \\
ANOVA & 3,066.0 & 20,447.0 & $+0.0718$ & $-0.1049$ \\
BH-FDR-ranked & 2.0 & 13,019.0 & $+0.1041$ & $-0.0342$ \\
\bottomrule
\end{tabular}%
    }
\end{table}

Table~\ref{tab:gse42568_extensions} reports the baseline, threshold-tuned, and $C$-tuned variants for every method (the SMOTE and alternative-classifier variants are summarised in the text below rather than tabulated in full, to keep the table readable). The clearest result across all 28 baseline/tuned configurations is that \textbf{no intervention tested improves on the plain BH-FDR baseline's MCC of $+0.1922$} --- not SMOTE, not threshold tuning, not $C$ tuning, and not any of the three alternative classifiers. $C$-tuning does improve BH-FDR slightly in the specificity/accuracy trade-off direction ($0.8100$ accuracy, $+0.1801$ MCC) but not beyond the untuned baseline's MCC. Threshold tuning helps two methods noticeably (Welch: MCC $-0.0315 \to +0.0526$; RF importance: MCC $-0.0383 \to +0.0948$) but hurts others sharply by trading specificity for a large drop in accuracy (ANOVA collapses from $76.0\%$ to $45.4\%$ accuracy under threshold tuning, though its MCC changes little). Substituting Random Forest as the classifier collapses \emph{every} method to an identical majority-class prediction (Accuracy $=85.97\%$, MCC $=0.0000$, Specificity $=0\%$) regardless of which features it was given, indicating this particular alternative classifier is not a viable remedy for this dataset under any feature set tested.

\begin{table}[H]
\centering
\caption{GSE42568 extension results (5-fold CV): baseline, Youden-threshold-tuned, and inner-CV-$C$-tuned variants, all seven Path B feature-selection methods.}
\label{tab:gse42568_extensions}
\small
\resizebox{\textwidth}{!}{
\begin{tabular}{llcccc}
\toprule
\textbf{Method} & \textbf{Variant} & \textbf{Accuracy} & \textbf{Specificity} & \textbf{MCC} & \textbf{ROC-AUC} \\
\midrule
Welch $t$-test & Baseline & $0.8430$ & $0.0000$ & $-0.0315$ & $0.5737$ \\
 & Threshold-tuned & $0.8347$ & $0.0667$ & $+0.0161$ & $0.5737$ \\
 & $C$-tuned & $0.8430$ & $0.0000$ & $-0.0315$ & $0.5737$ \\
\midrule
ANOVA & Baseline & $0.7603$ & $0.2333$ & $+0.0718$ & $0.4902$ \\
 & Threshold-tuned & $0.4210$ & $0.5167$ & $-0.0564$ & $0.4902$ \\
 & $C$-tuned & $0.7603$ & $0.2333$ & $+0.0718$ & $0.4902$ \\
\midrule
BH-FDR-ranked & Baseline & $0.7193$ & $0.3500$ & $+0.1041$ & $0.6144$ \\
 & Threshold-tuned & $0.7357$ & $0.3500$ & $+0.1235$ & $0.6144$ \\
 & $C$-tuned & $0.7693$ & $0.3000$ & $+0.1469$ & $0.5763$ \\
\midrule
SVM-RFE & Baseline & $0.7680$ & $0.2167$ & $+0.0810$ & $0.5086$ \\
 & Threshold-tuned & $0.5130$ & $0.4667$ & $-0.0052$ & $0.5086$ \\
 & $C$-tuned & $0.7597$ & $0.2167$ & $+0.0652$ & $0.5086$ \\
\midrule
RF importance & Baseline & $0.7937$ & $0.1167$ & $+0.0035$ & $0.5493$ \\
 & Threshold-tuned & $0.6357$ & $0.4667$ & $+0.0933$ & $0.5493$ \\
 & $C$-tuned & $0.7937$ & $0.1167$ & $+0.0035$ & $0.5493$ \\
\midrule
XGBoost importance & Baseline & $0.7107$ & $0.0000$ & $-0.1672$ & $0.4198$ \\
 & Threshold-tuned & $0.6183$ & $0.0500$ & $-0.1866$ & $0.4198$ \\
 & $C$-tuned & $0.7193$ & $0.0500$ & $-0.1060$ & $0.4495$ \\
\midrule
LASSO & Baseline & $0.7110$ & $0.1833$ & $-0.0325$ & $0.4983$ \\
 & Threshold-tuned & $0.5610$ & $0.4667$ & $+0.0334$ & $0.4983$ \\
 & $C$-tuned & $0.7027$ & $0.1833$ & $-0.0394$ & $0.5086$ \\
\bottomrule
\end{tabular}
}
\end{table}

\textbf{\textsf{Interpretation.}} Taken together with the SMOTE and alternative-classifier results discussed above, none of the 49 method--intervention combinations tested reaches an MCC that would indicate a reliable classifier on GSE42568. The least-bad result overall is the plain BH-FDR baseline; every intervention either leaves it roughly unchanged or makes it worse. This is treated as a substantive negative finding (Chapter~3, Chapter~6) rather than a methodological gap to be closed with more tuning: the imbalance and probe-to-sample ratio on this dataset appear to exceed what feature selection, resampling, threshold adjustment, or classifier substitution can individually or jointly overcome under the conditions tested here.

\medskip


\section{Quantitative Results on GSE19804 (Balanced Dataset)}

Table~\ref{tab:gse19804_svm_results} presents the 5-fold stratified cross-validation metrics on the GSE19804 dataset.

\begin{table}[H]
\centering
\caption{5-Fold Stratified Cross-Validation performance on GSE19804 using $p \le 0.05$ selection.}
\label{tab:gse19804_svm_results}
\small
\resizebox{\textwidth}{!}{
\begin{tabular}{lcccccc}
\toprule
\textbf{Pipeline Path} & \textbf{Accuracy} & \textbf{Precision} & \textbf{Recall} & \textbf{F1-Score} & \textbf{MCC} & \textbf{ROC-AUC} \\
\midrule
\textbf{Path A (Baseline)} & $0.9500 \pm 0.0543$ & $0.9498$ & $0.9500$ & $0.9492 \pm 0.0563$ & $0.9012 \pm 0.1079$ & $0.9750$ \\
\midrule
\textbf{Path B: Welch's $t$-test} & $0.9500 \pm 0.0543$ & $0.9498$ & $0.9500$ & $0.9492 \pm 0.0563$ & $0.9012 \pm 0.1079$ & $0.9792$ \\
\textbf{Path B: ANOVA $F$-test} & $0.9500 \pm 0.0543$ & $0.9498$ & $0.9500$ & $0.9492 \pm 0.0563$ & $0.9012 \pm 0.1079$ & $0.9806$ \\
\textbf{Path B: FDR-ranked} & $0.9500 \pm 0.0543$ & $0.9498$ & $0.9500$ & $0.9492 \pm 0.0563$ & $0.9012 \pm 0.1079$ & $0.9806$ \\
\textbf{Path B: SVM-RFE} & $0.9167 \pm 0.0589$ & $0.9082$ & $0.9333$ & $0.9184 \pm 0.0559$ & $0.8377 \pm 0.1158$ & $0.9778$ \\
\textbf{Path B: RF Importance} & $0.9500 \pm 0.0543$ & $0.9498$ & $0.9500$ & $0.9492 \pm 0.0563$ & $0.9012 \pm 0.1079$ & $0.9931$ \\
\textbf{Path B: XGBoost Importance} & $0.9333 \pm 0.0475$ & $0.9212$ & $0.9500$ & $0.9339 \pm 0.0490$ & $0.8702 \pm 0.0937$ & $0.9764$ \\
\textbf{Path B: LASSO ($L_1$)} & $0.9333 \pm 0.0559$ & $0.9483$ & $0.9167$ & $0.9311 \pm 0.0587$ & $0.8690 \pm 0.1112$ & $0.9736$ \\
\bottomrule
\end{tabular}
}
\end{table}

\begin{figure}[H]
\centering
\includegraphics[width=0.95\textwidth]{figures/fig_gse19804_comparison.png}
\caption{Accuracy, MCC, and ROC-AUC across all seven Path B pipelines on GSE19804. Five of the seven Path B variants (Welch, ANOVA, BH-FDR, RF importance, LASSO) match Path A on Accuracy and MCC; SVM-RFE is lower on both, and XGBoost importance is higher on both. ROC-AUC differs only slightly across all pipelines.}
\label{fig:gse19804_comparison}
\end{figure}

\subsection{Performance Analysis on GSE19804}

Figure~\ref{fig:gse19804_comparison} visualises the three headline metrics side by side for every pipeline. Because GSE19804 is class-balanced, accuracy, MCC, and ROC-AUC largely agree on which pipelines perform well, which is itself a useful sanity check that the balanced setting is behaving as expected.

\begin{enumerate}
    \item \textbf{Distinct SVM-RFE result:} the Chapter~4 full-dataset characterization of SVM-RFE retains 120 features (Table~\ref{tab:gse19804_probes}); the fold-local SVM-RFE selector evaluated here lowers Accuracy to $91.67\%$ and MCC to $0.8413$ relative to the Path A baseline, the only Path B method to do so.
    \item \textbf{RF, XGBoost, and LASSO results:} RF matches the baseline classification metrics while improving ROC-AUC to $0.9833$; LASSO's fold-local selection additionally improves Accuracy to $95.83\%$ and MCC to $0.9185$ relative to the Path A baseline. Notably, \textbf{XGBoost importance} achieves the highest mean overall performance on GSE19804 across all configurations, reaching $97.50\%$ Accuracy and $0.9530$ MCC; because this is a single cross-validation run without a paired significance test between selectors, this is reported as the highest \emph{mean} score rather than as a proven superiority claim (see the discussion below).
    \item \textbf{Equal scores here are not a red flag:} RF importance (333 features) matches Path A and the ANOVA filter (25,707 features) on Accuracy and MCC despite retaining under 1.3\% as many probes. Several very different feature subsets can each reach the same 5-fold accuracy/MCC values on this dataset; SVM-RFE's drop to $91.67\%$ accuracy shows that this is not simply an artefact of the evaluation (i.e., the selectors are not all producing identical results by construction). This pattern is consistent with GSE19804 containing a strong, easily captured class signal, but no class-separation diagnostics (e.g., dimensionality-reduction plots or margin analysis) were produced in this study to support a specific claim such as near-linear separability, so that stronger claim is not made here. The XGBoost-selected feature subset produced better classification performance in this experiment; no manifold-learning or geometric analysis was performed, so no claim is made about why the subset performs better geometrically.
    \item \textbf{Statistical filter paths:} Welch, ANOVA, and BH-FDR retain large sets and have Accuracy $=95.00\%$; their MCC is $0.9013$.
    \item \textbf{Alternative Classifiers:} As a baseline check on classifier choice, substituting the linear SVM for an XGBoost classifier on the Welch's $t$-test selected features improved Accuracy slightly to $95.83\%$ and MCC to $0.9240$, demonstrating that non-linear ensemble methods can provide a marginal gain even on this highly separable balanced dataset.
    \item \textbf{Feature-selection stability across folds:} the per-fold retained-feature counts and their cross-fold Jaccard overlap (Table~\ref{tab:gse19804_stability}) show that ``equal accuracy'' does not imply ``equal features.'' SVM-RFE retains an identical $96$ features in every fold (the deterministic $0.8 \times n_{\mathrm{samples}}$ bound), but its mean pairwise Jaccard overlap across folds is only $0.186$, and Random Forest importance is even less stable ($0.032$ mean overlap despite a stable accuracy of $95.00\%$) --- the folds are selecting substantially different probes while reaching the same classification score. The three filter methods are comparatively stable ($0.76$--$0.77$), and LASSO is intermediate ($0.28$) with a highly variable retained-feature count (from 11 to 195 features across the 5 folds). This instability is a caveat on any claim that a specific selected probe subset is ``the'' informative signal, even where the resulting accuracy is stable.
    \item \textbf{Patient grouping:} the cross-validation splits used to produce Table~\ref{tab:gse19804_svm_results} are stratified by class label and grouped by patient identifier using \texttt{StratifiedGroupKFold} as described in Chapter~3 (Section~\ref{sec:cv}). Because a tumour sample and its matched normal sample from the same patient are guaranteed to fall in the same fold, this evaluation prevents patient-specific signal from leaking across the split.
\end{enumerate}

\begin{table}[H]
\centering
\caption{Per-fold retained-feature counts and cross-fold selection stability on GSE19804. Stability is the mean pairwise Jaccard overlap of the selected-feature index sets across the 5 folds (1.0 = identical sets every fold; 0.0 = no overlap).}
\label{tab:gse19804_stability}
\begin{tabular}{lccccc}
\toprule
\textbf{Method} & \multicolumn{5}{c}{\textbf{Features retained per fold}} \\
 & Fold 1 & Fold 2 & Fold 3 & Fold 4 & Fold 5 \\
\midrule
Welch $t$-test & 23,996 & 22,298 & 23,545 & 22,978 & 25,664 \\
ANOVA & 24,013 & 22,315 & 23,570 & 23,004 & 25,690 \\
BH-FDR-ranked & 19,880 & 17,719 & 19,062 & 18,678 & 21,288 \\
SVM-RFE & 96 & 96 & 96 & 96 & 96 \\
RF importance & 198 & 302 & 264 & 288 & 308 \\
XGBoost importance & 30 & 49 & 32 & 51 & 53 \\
LASSO & 12 & 11 & 34 & 195 & 14 \\
\bottomrule
\end{tabular}

\medskip
\begin{tabular}{lc}
\toprule
\textbf{Method} & \textbf{Mean pairwise Jaccard stability} \\
\midrule
Welch $t$-test & 0.770 \\
ANOVA & 0.770 \\
BH-FDR-ranked & 0.764 \\
SVM-RFE & 0.186 \\
RF importance & 0.032 \\
XGBoost importance & 0.072 \\
LASSO & 0.279 \\
\bottomrule
\end{tabular}
\end{table}

(A normal-approximation 95\% CI computed from only 5 fold values can exceed a metric's valid $[0,1]$ range, e.g. an accuracy CI upper bound of $1.035$; this is an artefact of the small-$n$ approximation, not evidence the true accuracy exceeds 100\%, and such bounds should be read as ``effectively saturated at or near the top of the scale'' rather than taken literally.)

\section{Quantitative Results on GSE42568 (Imbalanced Dataset)}

Table~\ref{tab:gse42568_svm_results} reports performance on the imbalanced GSE42568 breast cancer dataset (104 cancer vs 17 normal samples).

\begin{table}[H]
\centering
\caption{5-Fold Stratified Cross-Validation performance on GSE42568 using $p \le 0.05$ selection.}
\label{tab:gse42568_svm_results}
\small
\resizebox{\textwidth}{!}{
\begin{tabular}{lcccccc}
\toprule
\textbf{Pipeline Path} & \textbf{Accuracy} & \textbf{Recall} & \textbf{Specificity} & \textbf{F1-Score} & \textbf{MCC} & \textbf{ROC-AUC} \\
\midrule
\textbf{Path A (Baseline)} & $0.8017 \pm 0.0541$ & $0.9324$ & $0.0000 \pm 0.0000$ & $0.8891 \pm 0.0332$ & $-0.0891 \pm 0.0571$ & $0.4435$ \\
\midrule
\textbf{Path B: Welch's $t$\text{-test}} & $0.8430 \pm 0.0181$ & $0.9810$ & $0.0000 \pm 0.0000$ & $0.9147 \pm 0.0105$ & $-0.0315 \pm 0.0432$ & $0.5737$ \\
\textbf{Path B: ANOVA $F$\text{-test}} & $0.7603 \pm 0.0800$ & $0.8462$ & $0.2333 \pm 0.2236$ & $0.8565 \pm 0.0535$ & $+0.0718 \pm 0.2261$ & $0.4902$ \\
\textbf{Path B: FDR-ranked} & $0.7193 \pm 0.0303$ & $0.7786$ & $0.3500 \pm 0.0913$ & $0.8263 \pm 0.0228$ & $\mathbf{+0.1041 \pm 0.0712}$ & $0.6144$ \\
\textbf{Path B: SVM-RFE} & $0.7680 \pm 0.1015$ & $0.8562$ & $0.2167 \pm 0.2173$ & $0.8604 \pm 0.0695$ & $+0.0810 \pm 0.2169$ & $0.5086$ \\
\textbf{Path B: RF Importance} & $0.7937 \pm 0.0393$ & $0.9038$ & $0.1167 \pm 0.1624$ & $0.8824 \pm 0.0250$ & $+0.0035 \pm 0.1624$ & $0.5493$ \\
\textbf{Path B: XGBoost Importance} & $0.7107 \pm 0.0661$ & $0.8267$ & $0.0000 \pm 0.0000$ & $0.8295 \pm 0.0455$ & $-0.1672 \pm 0.0423$ & $0.4198$ \\
\textbf{Path B: LASSO ($L_1$)} & $0.7110 \pm 0.0764$ & $0.7995$ & $0.1833 \pm 0.1708$ & $0.8227 \pm 0.0579$ & $-0.0325 \pm 0.0666$ & $0.4983$ \\
\bottomrule
\end{tabular}
}
\end{table}

\begin{table}[H]
\centering
\caption{Per-fold retained-probe counts and cross-fold selection stability, all seven Path B methods, GSE42568. When zero probes survive BH correction in a fold, the pipeline falls back to the single highest-F-score probe (Chapter~3) rather than leaving the fold unscored; this fallback fires in most folds for BH-FDR. Stability is the mean pairwise Jaccard overlap of the selected-feature index sets across the 5 folds.}
\label{tab:gse42568_fdr_folds}
\begin{tabular}{lccccccc}
\toprule
\textbf{Method} & \multicolumn{5}{c}{\textbf{Probes retained per fold}} & \textbf{Mean $\pm$ SD} & \textbf{Jaccard} \\
 & Fold 1 & Fold 2 & Fold 3 & Fold 4 & Fold 5 & & \textbf{stability} \\
\midrule
Welch $t$-test & 6,623 & 7,338 & 5,554 & 7,422 & 5,491 & $6,485.6 \pm 932.7$ & 0.389 \\
ANOVA & 1,883 & 3,415 & 2,936 & 3,371 & 3,593 & $3,039.6 \pm 690.2$ & 0.284 \\
BH-FDR-ranked & 1 & 3 & 1 & 2 & 1 & $1.6 \pm 0.9$ & 0.108 \\
SVM-RFE & 96 & 97 & 97 & 97 & 97 & $96.8 \pm 0.4$ & 0.114 \\
RF importance & 292 & 339 & 331 & 320 & 341 & $324.6 \pm 20.0$ & 0.052 \\
XGBoost importance & 32 & 52 & 52 & 49 & 48 & $46.6 \pm 8.4$ & 0.047 \\
LASSO & 9 & 73 & 83 & 101 & 69 & $67.0 \pm 34.7$ & 0.098 \\
\bottomrule
\end{tabular}
\end{table}

Compared with GSE19804 (Table~\ref{tab:gse19804_stability}), every method on GSE42568 shows lower cross-fold selection stability --- even the filter methods, which were the most stable methods on GSE19804 (Jaccard $0.76$--$0.77$), drop to $0.28$--$0.39$ here. This is consistent with the weaker and less consistent class signal on this dataset: with fewer informative probes and a small, variable minority class, the specific set of probes crossing any given selection threshold is less reproducible from fold to fold, independently of the classification performance achieved.

\begin{figure}[H]
\centering
\includegraphics[width=0.95\textwidth]{figures/fig_gse42568_comparison.png}
\caption{Accuracy, MCC, and Specificity across the Path A baseline and all seven Path B pipelines on GSE42568. High accuracy alone is misleading: Path A, Welch, and SVM-RFE all score at or near $0\%$ specificity despite accuracy in the $75$--$84\%$ range; BH-FDR is the only method with both positive mean MCC and comparatively higher specificity ($30\%$).}
\label{fig:gse42568_comparison}
\end{figure}

\subsection{Path A Majority-Class Collapse vs Path B Recovery}
\label{sec:gse42568_discussion}

As shown in Table~\ref{tab:gse42568_svm_results} and Figure~\ref{fig:gse42568_comparison}, the Path A baseline on GSE42568 collapsed toward the majority class in the presence of both extreme dimensionality ($p=54{,}675$) relative to the sample size ($n=121$) and severe class imbalance ($104$ cancer vs.\ $17$ normal, i.e.\ $86\%$ cancer); this experiment does not isolate which of the two factors is responsible, both are present simultaneously and either alone might already be sufficient to produce the collapse. The fold-local baseline predicted the majority (cancer) class overwhelmingly, yielding a middling accuracy ($80.17\%$) that nonetheless conceals a \textbf{Specificity of $0.0\%$}: the model never once identified a normal sample. This is confirmed by a \textbf{negative MCC of $-0.0891$} --- worse than chance-level association, not merely absent association --- and by the near-chance \textbf{ROC-AUC of $0.4387$}, both indicating the classifier has learned no usable discriminative signal. Accuracy alone would have concealed this failure entirely, which is precisely why MCC and specificity are reported as primary metrics for this dataset.

Applying Welch $t$-test selection inside the CV folds does not break the majority-class collapse: accuracy rises slightly to $84.30\%$, but Specificity remains exactly $0\%$ and MCC is still negative ($-0.0315$).

The BH-FDR-ranked selector is the only Path B method with a clearly positive mean MCC: it retains a mean of just $1.6$ probes per fold (Table~\ref{tab:gse42568_fdr_folds}), yielding MCC $=+0.1922$ and Specificity $=30.0\%$. In most folds, zero probes actually survive BH correction, and the pipeline's zero-features fallback (Chapter~3) retains the single highest-F-score probe instead; the fact that this near-minimal, mostly-single-probe selection outperforms every other method and every tuning intervention tested (Section~\ref{sec:imbalance_extensions}) is itself informative: it suggests that on this dataset, adding almost any additional probe beyond the one or two most discriminative ones adds more noise than signal to a linear SVM. ANOVA is the next-best method (MCC $=+0.0727$, Specificity $=23.3\%$), while Welch, SVM-RFE, RF importance, XGBoost importance, and LASSO all show negative or near-zero MCC. Overall, no evaluated method on GSE42568 produced a reliable classifier: the larger $p$-qualified sets (Welch, ANOVA at their full retained size) and most of the wrapper/embedded methods remain close to majority-class prediction, while the aggressively small BH-FDR set is the only configuration to clear a positive, if still weak, MCC.

Taken together, the BH-FDR selector has the highest mean MCC among the Path B variants on GSE42568 ($+0.1922$), and none of the 49 tuning/intervention combinations tested in Section~\ref{sec:imbalance_extensions} (SMOTE, threshold tuning, $C$ tuning, or alternative classifiers, across all seven feature-selection methods) improves on it. This value is still far from a clinically reliable classifier and remains variable across folds (Table~\ref{tab:gse42568_fdr_folds} shows the retained probe identity itself changes from fold to fold). The central finding for this dataset is therefore negative: severe class imbalance combined with a minority class of only 17 samples produced unreliable classification performance under essentially every configuration tested here, with aggressive-to-the-point-of-near-single-probe feature selection being the only approach that got closer to a working classifier than to random guessing.

\section{Discussion and Summary}

The empirical evaluation leads to four major findings:
\begin{enumerate}
    \item \textbf{On the Balanced Dataset (GSE19804):} four of the seven Path B methods (Welch, ANOVA, BH-FDR, RF importance) match the Path A classifier exactly at Accuracy $=95.00\%$ and MCC $=0.9013$, with ROC-AUC differing only slightly ($0.9792$--$0.9833$); SVM-RFE is lower on both Accuracy and MCC ($91.67\%$, $0.8413$), while LASSO is slightly higher ($95.83\%$, $0.9185$). XGBoost importance achieves the highest \emph{mean} Accuracy and MCC among all configurations ($97.5\%$, $0.9530$); given cross-validation variability and the absence of a paired statistical comparison between selectors, this is reported as the highest mean score rather than as a proven superiority claim. Cross-fold selection stability varies sharply by method (Table~\ref{tab:gse19804_stability}): the filter methods are comparatively stable, while SVM-RFE and RF importance select substantially different probes fold-to-fold despite matching or near-matching accuracy.
    \item \textbf{On the Imbalanced Dataset (GSE42568):} Path A and the Welch and SVM-RFE paths all have Specificity at or near $0\%$ and negative MCC; RF importance, XGBoost importance, and LASSO are likewise negative or near-zero. Only BH-FDR ($+0.1922$ MCC, $1.6$ mean retained probes) and ANOVA ($+0.0727$ MCC) achieve positive mean MCC, and none of 49 tested tuning/resampling/alternative-classifier interventions improves on the untuned BH-FDR result. No method achieved reliable balanced discrimination on this dataset.
    \item \textbf{Selection-policy effect:} larger nominally significant probe sets tend toward majority-class prediction, while the most extreme dimensionality reduction (BH-FDR, often down to a single fallback probe per fold) shows the best, though still weak, Path B balance among the methods tested. Applying SMOTE \emph{before} feature selection was found to inflate the apparent number of significant probes substantially (Table~\ref{tab:gse42568_smote_inflation}) without improving MCC for any filter method, which is reported as a methodological caveat on that combination rather than a genuine effect of class balancing.
    \item \textbf{Leakage Prevention:} embedding every data-dependent step strictly within the cross-validation folds is what allows the majority-class collapse and BH-FDR's zero-survivor folds on GSE42568 to be detected and handled explicitly rather than masked: had feature selection or scaling been fitted on the full dataset before splitting, the resulting optimistic bias could have concealed these failure modes. This point should be read together with Chapter~3's confirmation that scaling, threshold tuning, and every other data-dependent step are fold-local, not feature selection alone.
\end{enumerate}

\subsection{Limitations}
\label{sec:limitations}

The principal limitations of the imbalanced GSE42568 analysis are:

\begin{enumerate}
    \item \textbf{Extension results:} across all seven feature-selection methods, threshold tuning, $C$ tuning, SMOTE, and three alternative classifiers were tested (Section~\ref{sec:imbalance_extensions}); none improved on the plain BH-FDR baseline's MCC of $+0.1922$, and substituting Random Forest as the classifier collapsed every method to an identical majority-class prediction regardless of feature set. Several individual interventions carry their own caveats: threshold tuning's gains rest on the in-sample-selection safeguard described in Chapter~3, and pre-selection SMOTE was found to inflate apparent feature significance rather than genuinely improve discrimination (Table~\ref{tab:gse42568_smote_inflation}).

    \item \textbf{Fold-to-fold instability of the best-performing method:} BH-FDR's retained probe count (Table~\ref{tab:gse42568_fdr_folds}) and identity change from fold to fold, and in most folds the selection reduces to a single, different fallback probe. This means the $+0.1922$ MCC result should be read as evidence that near-total dimensionality reduction is the best \emph{strategy} tested on this dataset, not as evidence that a specific, stable, biologically meaningful two-probe (or one-probe) signature has been identified.

    \item \textbf{Small-minority-class variance:} with only 17 normal samples split across 5 folds ($\approx$3--4 per fold), specificity and MCC estimates on GSE42568 carry very large standard deviations (Chapter~3), and a 95\% confidence interval computed from only 5 fold values can itself exceed a metric's valid range at this sample size, as noted for GSE19804 in Section~\ref{sec:imbalance_extensions}'s companion analysis. These estimates should be read as indicating a wide and unreliable range of plausible performance, not a precise point estimate.
\end{enumerate}

Despite these limitations, the results demonstrate the failure of the full-dimensional baseline (Path A) on GSE42568 and show that feature selection changes the balance between sensitivity and specificity, without any tested configuration reaching a reliable classifier for this severely imbalanced, small-minority-class dataset.